\documentclass[tikz,border=4pt]{standalone}

\usepackage{amsmath}
\usetikzlibrary{arrows.meta}

\definecolor{rodblue}{rgb}{0.02,0.16,0.58}
\definecolor{anglered}{rgb}{0.62,0.00,0.02}

\begin{document}

\begin{tikzpicture}[
    x=1cm,
    y=1cm,
    line cap=round,
    line join=round,
    force/.style={
        draw=black,
        line width=2.8pt,
        -{Stealth[length=4.2mm,width=3.2mm]}
    },
    extension/.style={
        draw=black,
        line width=1.9pt,
        dash pattern=on 6pt off 5pt
    },
    normal/.style={
        draw=anglered,
        line width=1.9pt,
        dash pattern=on 6pt off 5pt
    },
    angle arc/.style={
        draw=anglered,
        line width=2.2pt
    },
    force label/.style={
        text=black,
        font=\fontsize{24}{28}\selectfont
    },
    angle label/.style={
        text=anglered,
        font=\fontsize{22}{26}\selectfont
    }
]

% ============================================================
% 1. 关键点
% ============================================================

\coordinate (A) at (0,0);
\coordinate (B) at (2.15,3.05);
\coordinate (C) at (1.204,1.708);

% 蓝色杆件方向单位向量约为：
% u = (0.5762, 0.8173)

\coordinate (LowerExtension) at (-0.749,-1.063);
\coordinate (UpperExtension) at (2.830,4.014);

% T_5、T_6 端点
\coordinate (TfiveEnd) at (2.558,4.700);
\coordinate (TsixEnd)  at (-1.795,-0.761);

% ============================================================
% 2. 蓝色杆件
% ============================================================

\draw[
    draw=rodblue,
    line width=4.2pt
]
(A) -- (B);

% ============================================================
% 3. 杆件延长线
% ============================================================

\draw[extension]
(A) -- (LowerExtension);

\draw[extension]
(B) -- (UpperExtension);

% ============================================================
% 4. 顶部竖直法线
% ============================================================

\draw[normal]
(B) -- ++(0,1.28);

% ============================================================
% 5. A 点受力
% ============================================================

% F_s
\draw[force]
(A) -- ++(0,1.18);

% G_s
\draw[force]
(A) -- ++(0,-1.72);

% T_6
\draw[force]
(A) -- (TsixEnd);

% ============================================================
% 6. C 点受力
% ============================================================

% F_5
\draw[force]
(C) -- ++(0,1.45);

% G_5
\draw[force]
(C) -- ++(0,-2.15);

% ============================================================
% 7. B 点受力
% ============================================================

% T_5
\draw[force]
(B) -- (TfiveEnd);

% ============================================================
% 8. 作用点
% ============================================================

\fill[black] (A) circle[radius=2.1pt];
\fill[black] (B) circle[radius=2.1pt];
\fill[black] (C) circle[radius=1.65pt];

% ============================================================
% 9. 力的标注
% ============================================================

\node[
    force label,
    anchor=south east
] at (-0.15,1.32) {$F_s$};

\node[
    force label,
    anchor=north
] at (-0.07,-1.90) {$G_s$};

\node[
    force label,
    anchor=north east
] at (-1.895,-0.861) {$T_6$};

\node[
    force label,
    anchor=south east
] at (1.044,3.278) {$F_5$};

\node[
    force label,
    anchor=north west
] at (1.304,-0.542) {$G_5$};

\node[
    force label,
    anchor=south west
] at (2.608,4.740) {$T_5$};

% ============================================================
% 10. A 点的 theta_6
% T_6 与竖直向下方向的夹角
% ============================================================

\draw[angle arc]
($(A)+(-157.02:0.55)$)
arc[
    start angle=-157.02,
    end angle=-90,
    radius=0.55
];

\node[angle label]
at (-0.43,-0.52) {$\theta_6$};

% ============================================================
% 11. A 点下部的 varphi_5
% 杆件左下延长线与竖直向下方向的夹角
% ============================================================

\draw[angle arc]
($(A)+(-125.18:0.94)$)
arc[
    start angle=-125.18,
    end angle=-90,
    radius=0.94
];

\node[angle label]
at (-0.18,-0.94) {$\varphi_5$};

% ============================================================
% 12. C 点的 varphi_5
% 杆件左下方向与 G_5 的夹角
% ============================================================

\draw[angle arc]
($(C)+(-125.18:0.43)$)
arc[
    start angle=-125.18,
    end angle=-90,
    radius=0.43
];

\node[angle label]
at (0.934,1.218) {$\varphi_5$};

% ============================================================
% 13. B 点的 theta_5
% 顶部竖直法线与 T_5 的夹角
% ============================================================

\draw[angle arc]
($(B)+(90:0.39)$)
arc[
    start angle=90,
    end angle=76.10,
    radius=0.39
];

\node[angle label]
at (2.22,3.57) {$\theta_5$};

% ============================================================
% 14. B 点顶部的 varphi_5
% 杆件右上延长线与顶部法线的夹角
% ============================================================

\draw[angle arc]
($(B)+(54.82:0.72)$)
arc[
    start angle=54.82,
    end angle=90,
    radius=0.72
];

\node[angle label]
at (2.77,3.72) {$\varphi_5$};

% ============================================================
% 15. 绘图范围
% ============================================================

\path[use as bounding box]
(-2.65,-2.45) rectangle (4.35,5.15);

\end{tikzpicture}

\end{document}